\section{Derivaciones financieras y asignativas que sustentan el mecanismo}
\label{app:compact_financial_derivations}
Esta sección presenta cuatro argumentos completos del dossier previo,
ahora en forma autocontenida. Se conservan sus etiquetas originales para
hacer posible el cotejo con el registro de resultados y con la formalización.
Las hipótesis declaradas son suficientes para cada afirmación, pero no
convierten estas conclusiones en existencia de equilibrio dinámico.

\begin{lemma}[Suficiencia del colateral para repago]
\label{lem:repayment}
Suponga $0\leq\delta<1$, $0\leq\theta_i<1$,
$R^L,R^D,\Pi',q'>0$ y $a,b,k'\geq0$.
En cada estado admisible, los beneficios operativos cumplen
$\pi^{op}{}'\geq0$ y el contrato de garantía robusta implica
\[
\frac{R^L b}{\Pi'}\leq
\theta_i(1-\delta)q'k'.
\]
Si el patrimonio siguiente es
\[
N'=\pi^{op}{}'+(1-\delta)q'k'
-\frac{R^L b}{\Pi'}+\frac{R^D a}{\Pi'},
\]
el préstamo queda cubierto por el valor pignorable del capital
remanente y $N'\geq0$. El patrimonio es estrictamente positivo
cuando $k'>0$ o $a>0$.
\end{lemma}
\begin{proof}
La desigualdad contractual se ha formulado después de dividir
la promesa nominal por la inflación realizada positiva, de modo
que conserva el sentido de la desigualdad. Sustituyendo su cota
superior en el término negativo del balance de patrimonio,
\begin{align*}
N'&=\pi^{op}{}'+(1-\delta)q'k'
-\frac{R^L b}{\Pi'}+\frac{R^D a}{\Pi'}\\
&\geq
\pi^{op}{}'+(1-\theta_i)(1-\delta)q'k'
+\frac{R^D a}{\Pi'}\geq0.
\end{align*}
La última desigualdad sigue de las hipótesis de signos.
Si $k'>0$, entonces $(1-\theta_i)(1-\delta)q'k'>0$;
si $a>0$, entonces $R^D a/\Pi'>0$. No se demuestra
por ello solvencia agregada del banco ni consumo empresarial
positivo: ambas propiedades requieren restricciones adicionales.
\end{proof}

\begin{lemma}[Incompatibilidad de ahorro y endeudamiento interiores]
\label{lem:v0m_original_6_2}
Considere una empresa con $R^L>R^D>0$, descuento
$\beta_E>0$, inflación futura $\Pi'>0$ y
utilidad marginal de patrimonio $V'_N>0$ casi seguramente.
Suponga que las expectativas indicadas son finitas y que
las condiciones de primer orden de depósito y préstamo,
para posiciones interiores $a,b>0$, son
\[
\lambda=\beta_E\mathbb E_t[V'_N R^D/\Pi'],
\quad
\lambda=\beta_E\mathbb E_t[V'_N R^L/\Pi']+\mu R^L,
\quad \mu\geq0.
\]
Entonces ninguna solución óptima que satisfaga esas
condiciones puede tener simultáneamente $a>0$ y $b>0$.
\end{lemma}
\begin{proof}
Supongamos por contradicción que ambas posiciones son
estrictamente positivas, de modo que sus duales de no
negatividad son nulos. Al igualar las dos expresiones para
$\lambda$ y restar,
\[
0=\beta_E\mathbb E_t
\left[V'_N\frac{R^L-R^D}{\Pi'}\right]+\mu R^L.
\]
El primer término es estrictamente positivo por las
hipótesis de signos y el segundo no negativo.
Esto contradice la igualdad a cero. Las condiciones KKT
son aquí necesarias; no se ha demostrado suficiencia global
del problema intertemporal a partir de ellas.
\end{proof}

\begin{proposition}[Efecto parcial de abaratar el préstamo]
\label{prop:v0m_original_6_3}
Sea $n:=N-c-a>0$, $q>0$, $H>0$ y $R^L>0$,
con $q>H/R^L$. Suponga el presupuesto empresarial
$qk'=n+b$, el límite de colateral $R^L b\leq Hk'$
y $b,k'\geq0$. El máximo capital financiable es
\[
k^{\max}=\frac{n}{q-H/R^L},\qquad
\frac{\partial k^{\max}}{\partial R^L}
=-\frac{nH}{(R^L)^2(q-H/R^L)^2}<0.
\]
\end{proposition}
\begin{proof}
La garantía implica $b\leq Hk'/R^L$. Sustituyéndola
en el presupuesto se obtiene
$qk'\leq n+Hk'/R^L$. Como el denominador es
estrictamente positivo,
\[
0\leq k'\leq \frac{n}{q-H/R^L}.
\]
El máximo es alcanzable con $b=Hk'/R^L\geq0$,
que satisface el presupuesto por construcción.
Aplicando la regla de la cadena al recíproco,
\[
\frac{\partial}{\partial R^L}
\frac{n}{q-H/R^L}
=-\frac{n}{(q-H/R^L)^2}
\frac{H}{(R^L)^2}<0.
\]
La comparación mantiene fijos $n,q,H$.
Una mayor capacidad financiable no demuestra
un mayor capital óptimo ni una restricción menos
vinculante en equilibrio general.
\end{proof}

\begin{lemma}[Referencia eficiente y derivada de reasignación]
\label{lem:allocation_loss_identity}
Sean $v_i,A_i,k_i>0$ para $i=H,L$, $L>0$,
$\alpha,\beta>0$ y $d=1-\alpha-\beta>0$.
Manténganse fijas las masas, productividades y elasticidades.
En la economía con tecnología
$Y_i=A_i k_i^\alpha\ell_i^\beta$ y
trabajo total $\sum_i v_i\ell_i=L$, defina
\[
G=\sum_i v_i A_i^{1/(1-\beta)}
k_i^{\alpha/(1-\beta)},\quad
K=\sum_i v_i k_i,\quad Z_A=\sum_i v_iA_i^{1/d}.
\]
Con empleo óptimamente distribuido entre empresas,
\[
Y=L^\beta G^{1-\beta},\quad
Y^{eff}=K^\alpha L^\beta Z_A^d,\quad
\mathscr L=\log(Y^{eff}/Y)\geq0.
\]
Si $\kappa_H=v_Hk_H/K$ y
$\omega_H=v_HY_H/Y$, a lo largo de cualquier
trayectoria diferenciable positiva de capital,
\[
\mathrm d\mathscr L=
\alpha(\kappa_H-\omega_H)
\mathrm d\log(k_H/k_L).
\]
Finalmente, $\omega_H>\kappa_H$ si y solo si el
producto marginal físico de capital de H es
mayor que el de L, evaluados en el empleo óptimo.
\end{lemma}
\begin{proof}
Dado el capital, el problema cóncavo de empleo es
$\max_{\ell_i\geq0}\sum_i v_iA_i k_i^\alpha\ell_i^\beta$
sujeto a $\sum_i v_i\ell_i=L$.
Las condiciones de primer orden interiores igualan
$\beta A_i k_i^\alpha\ell_i^{\beta-1}$ entre empresas.
Escribiendo $g_i=A_i^{1/(1-\beta)}
k_i^{\alpha/(1-\beta)}$, resultan
$\ell_i=Lg_i/G$. En efecto,
$\sum_i v_i\ell_i=L(\sum_i v_i g_i)/G=L$,
y la estricta concavidad en trabajo certifica
que se trata del máximo. Al sustituir se obtiene
\[
Y=\sum_i v_iA_i k_i^\alpha (Lg_i/G)^\beta
=L^\beta G^{1-\beta}.
\]

Para permitir también reasignación de capital, tome
$k_i^*=K A_i^{1/d}/Z_A$ y
$\ell_i^*=L A_i^{1/d}/Z_A$. Satisfacen
$\sum_i v_i k_i^*=K$ y $\sum_i v_i\ell_i^*=L$.
En estos puntos tanto
$\alpha A_i(k_i^*)^{\alpha-1}(\ell_i^*)^\beta$
como
$\beta A_i(k_i^*)^\alpha(\ell_i^*)^{\beta-1}$
son independientes de $i$, porque sus factores
de $A_i$ tienen exponente total
$1+(\alpha+\beta-1)/d=0$.
La tecnología es cóncava para
$\alpha+\beta<1$; la desigualdad del plano de soporte,
sumada con pesos $v_i$, muestra que esta asignación
maximiza producto a dotaciones $(K,L)$.
Al sustituir los niveles óptimos resulta
$Y^{eff}=K^\alpha L^\beta Z_A^d$.
Por su optimalidad $Y^{eff}\geq Y>0$, y al tomar
logaritmos,
\[
\mathscr L=\alpha\log K+d\log Z_A
-(1-\beta)\log G\geq0.
\]

Con $\gamma:=\alpha/(1-\beta)$,
\begin{align*}
\mathrm d\log K&=\kappa_H\,\mathrm d\log k_H+
(1-\kappa_H)\,\mathrm d\log k_L,\\
\mathrm d\log G&=\gamma\big[
\omega_H\,\mathrm d\log k_H+
(1-\omega_H)\,\mathrm d\log k_L\big].
\end{align*}
La segunda igualdad usa
$\omega_i=v_i g_i/G=v_iY_i/Y$,
con el trabajo ya distribuido óptimamente.
Como $Z_A$ es fijo y $(1-\beta)\gamma=\alpha$,
restar las dos diferenciales produce
\[
\mathrm d\mathscr L=
\alpha(\kappa_H-\omega_H)
(\mathrm d\log k_H-\mathrm d\log k_L),
\]
como se afirmó; la variación de $L$ se cancela
incluso si el empleo total cambia a lo largo de
la trayectoria.

Finalmente, $MPK_i=\alpha Y_i/k_i$.
La desigualdad $\omega_H>\kappa_H$ equivale a
$\omega_H(1-\kappa_H)>
\kappa_H(1-\omega_H)$, y por las
definiciones de $\omega_H,\kappa_H$ esta última
equivale a $Y_H/k_H>Y_L/k_L$.
Multiplicar por $\alpha>0$ prueba
$MPK_H>MPK_L$. Por tanto, $A_H>A_L$
por sí sola no determina el signo de la
derivada asignativa.
\end{proof}

\section{Formalización en Lean: correspondencia de los ocho resultados}
\label{app:compact_lean_matrix}
El paper evaluable está fijado sobre una revisión específica;
la versión histórica de los otros 52 resultados se conserva
por separado en la rama
\texttt{archive/ie-paper-full-20261118}. La tabla no convierte
pruebas auxiliares o condicionales en equivalencias literales.
En la fecha del precierre, ninguna de las ocho filas tiene
una equivalencia integral paper--Lean acreditada en la matriz
conservadora; el estatus es \emph{abierto con evidencia parcial}.

\begingroup\small
\begin{longtable}{@{}p{1.05cm}p{4.0cm}p{3.1cm}p{5.8cm}@{}}
\toprule
Resultado & Enunciado del artículo & Prueba o módulo relacionado & Estado y limitación\\\midrule\endhead
\ref{prop:local_macro_multiplicity} &
Dos raíces locales de complementariedad condicionadas &
Conteo escalar (Lean) &
Abierto: el lema de conteo supone una raíz interior y
monotonía; falta puente con la función reducida $C^1$ y SOURCE.\\
\ref{prop:compact_exact_tail_power} &
Decaimiento por bloques del operador lineal decimal &
Certificado racional (Python) &
Abierto en Lean: cotas racionales comprobadas por Python;
la prueba de normas es condicional a ellas.\\
\ref{lem:compact_robust_tail_power} &
Robustez de matriz a $10^{-4}$ &
Certificado racional (Python) &
Abierto en Lean: majorantes racionales y telescopado
no están certificados como teorema literal en Lean.\\
\ref{prop:compact_source_collateral_effective_wedge} &
Cuña de precios efectivos bajo garantías activas &
Cuña de Euler (Lean) &
Abierto: el lema Euler--pignorabilidad verificado en SOURCE
no es la identidad literal de precios efectivos entre H y L.\\
\ref{lem:repayment} &
Repago robusto y patrimonio no negativo &
Interfaz de repago (Lean) &
Abierto en matriz: existe candidato kernel-verificado
en interfaz fuente; falta cotejar formalmente el paper
definitivo y su SHA.\\
\ref{lem:v0m_original_6_2} &
No ahorrar y endeudarse simultáneamente &
KKT financiera (Lean) &
Abierto en matriz: existe prueba relacionada sobre
KKT normalizadas; falta validar equivalencia de
hipótesis del manuscrito definitivo.\\
\ref{prop:v0m_original_6_3} &
Máximo financiable y derivada respecto de tasa &
Derivada financiera (Lean) &
Abierto en matriz: derivada auxiliar probada;
se requiere correspondencia literal de restricciones
y cuantificadores.\\
\ref{lem:allocation_loss_identity} &
Referencia eficiente y derivada de reasignación &
Asignación y MPK (Lean) &
Abierto en matriz: 42 declaraciones relacionadas;
faltan igualdad literal de la especificación y
cotejo integral de todas las hipótesis.\\
\bottomrule
\end{longtable}
\endgroup
Los nombres completos de las declaraciones y sus archivos se
encuentran en \path{paper/lean_formalization_matrix_8.md}; el
historial se conserva en \path{paper/lean_crosswalk_57.json}.
Ninguna ejecución previa constituye un chequeo nuevo de esta
revisión. El control oficial pendiente es:
\begin{quote}\small
\texttt{python3} \path{scripts/paper_contribution.py}
\texttt{check Ventura26LiquidityAllocation --fast}.
\end{quote}

\section{Colaboración con IA y veredictos documentados}
\label{app:compact_ai_log}
El autor reconoce asistencia en la formulación, las
derivaciones, la programación, la exposición y el examen
crítico. Los intercambios disponibles se preservan
en \texttt{prompts.md}; se prohíbe imputar respuestas
no recuperadas o aprobaciones personales inexistentes.
\begin{longtable}{@{}p{4.8cm}p{6.0cm}p{3.35cm}@{}}
\toprule
Afirmación o contribución & Revisión y contraste & Juicio atribuible al autor\\\midrule\endhead
$A_H>A_L$ implica $MPK_H>MPK_L$ &
Rechazada: la escala de capital y trabajo también
determina el producto marginal. Véase el Lema
\ref{lem:allocation_loss_identity}. &
El autor objetó explícitamente esa inferencia.\\
Una menor tasa garantiza mejor asignación &
Rechazada: la capacidad parcial y la respuesta
en equilibrio son conceptos distintos; las ramas
numéricas tienen signos opuestos. &
El autor pidió revisar las dos ramas; aceptación
de prueba general pendiente.\\
Dos raíces escalares prueban dos equilibrios dinámicos &
Rechazada: faltan recuperación SOURCE, optimalidad
y condiciones terminales; hay obstrucción con
antecedente heredado común. &
La presentación distingue ambos objetos;
veredicto formal del paper pendiente.\\
CI verde implica equivalencia literal de todas las pruebas &
Rechazada: la interfaz Lean conserva brechas
específicas y candidatos condicionales. &
El autor solicitó una auditoría completa;
no se imputa aprobación automática.\\
Guion oral como evidencia de demostración &
Rechazada: el guion apoya la exposición, pero
el apéndice a mano y la prueba formal son objetos
diferentes. &
El autor pidió un guion natural y señaló su
carácter inicialmente esquemático.\\
\bottomrule
\end{longtable}
Los cálculos manuscritos genuinos y el juicio crítico
final sobre cada prueba deben incorporarse en
\texttt{hand/}; estas tablas no son prueba de que
el autor ya haya escrito o validado sus hojas.

\section{Reproducción, límites y procedencia}
\label{app:compact_reproduction}
Las figuras, tablas y experimentos numéricos son
reproducibles, bajo los parámetros ilustrativos y
tolerancias del código, mediante
\texttt{python code/verify.py} desde la raíz del repositorio
tras instalar \texttt{code/requirements.txt}.
Las cuatro figuras/PDF oficiales se compilan
por el workflow \texttt{.github/workflows/build.yml}.
El cuerpo compacto no demuestra una respuesta
empírica para Perú, selección de rama, equilibrio
competitivo de horizonte infinito ni cambio
permanente de la tasa de crecimiento.
